\documentclass[10pt,a4paper]{amsart}
\usepackage[margin=19mm]{geometry}
\usepackage{amsmath,amssymb,mathtools}
\usepackage[scr=rsfs,cal=euler]{mathalfa}
\usepackage{xfrac}
\usepackage[unicode,hidelinks,bookmarks=false]{hyperref}
\newcommand{\ie}{i.e.\ }
\newcommand{\cf}{cf.\ }
\newcommand{\resp}{respectively\ }
\newcommand{\Subsection}{\subsection}
\providecommand{\nobreakdash}{\nobreak\hbox{-}}
\setlength{\emergencystretch}{2em}
\multlinegap0pt
\usepackage{bm}
\usepackage{bbm}
\usepackage{dsfont}
\usepackage{xspace}
\usepackage{cancel}
\usepackage{mathtools}
\usepackage{amsthm}
\makeatletter
\@ifundefined{theorem}{\newtheorem{theorem}{Theorem}}{}
\@ifundefined{definition}{\newtheorem{definition}[theorem]{Definition}}{}
\@ifundefined{proposition}{\newtheorem{proposition}[theorem]{Proposition}}{}
\makeatother
\newcommand{\DG}{\overline{\mathsf{D}}}
\newcommand{\npeh}{^{n+1/2}}
\newcommand{\npe}{^{n+1}}
\newcommand{\n}{^{n}}
\renewcommand{\vec}[1]{\ensuremath{\mbox{$\mathbf{#1}$}}}
\title[Energy-consistent integration of mechanical systems based on]{Energy-consistent integration of mechanical systems based on Livens principle}
\author{Philipp L. Kinon, Peter Betsch}
\date{Source: arXiv:2312.02825v3 (CC BY 4.0)}
\begin{document}
\maketitle

\noindent\emph{Source and licence.} Energy-consistent integration of mechanical systems based on Livens principle. Version arXiv:2312.02825v3 (CC BY 4.0), distributed under the Creative Commons Attribution 4.0 International licence, as recorded in the arXiv licence metadata for this exact version and shown on the abstract page https://arxiv.org/abs/2312.02825v3. Full rights notice: licenses/arxiv-2312.02825-CC-BY-4.0.txt.\par\smallskip

\noindent\textit{Editorial scope.} Counted unit: the Proposition whose source label is \texttt{None} in the article (source0.tex lines 404--406, source label \texttt{None}), delivered together with the article's own complete original proof of it (source0.tex lines 408--424, one \texttt{proof} environment of 17 lines). No mathematics is written, repaired or re-derived: statement and proof are the article's own text line for line. Required context retained uncounted: lagrangefunction (lines 169--175); generalized\_energy (lines 245--251); HP\_EL\_discrete (lines 368--376); HP\_EL1\_discrete (lines 370--375); HP\_EL2\_discrete (lines 370--375); HP\_EL3\_discrete (lines 370--375); HP\_EL4\_discrete (lines 370--375); discrete\_energy\_function (lines 396--401). Every definition, display and label of those lines is retained unchanged. Omitted from this package and not redistributed: 1--168, 176--244, 252--367, 376--395, 402--403, 407--407, 425--731. Nothing inside a retained excerpt is omitted or altered: every retained body is delivered exactly as the article's lines are written, so no omission receipt of any kind accompanies this package. Citations: the retained excerpts carry no citation command at all, so no bibliography entry of the article is transcribed and no reference is redistributed. Reference adaptation: none was needed and none was made; every reference inside the delivered unit resolves inside it, so no unresolved reference or citation is produced. Printed slips kept literal: no printed word, symbol, hypothesis or number of the counted statement or of its proof was altered. Layout and class: the article's own class is \texttt{multibody2023\_full\_paper} with packages including import, color, amsthm, pgfplots, pgfplotstable, tikz, url; the wrapper is the lane's pinned \texttt{amsart} editorial wrapper at 10pt on A4, which restates the theorem-like environments the retained text needs and adds the article's own macro definitions that the retained text needs (\texttt{\textbackslash DG}, \texttt{\textbackslash n}, \texttt{\textbackslash npe}, \texttt{\textbackslash npeh}, \texttt{\textbackslash vec}), with extra wrapper packages (bm, bbm, dsfont, xspace, cancel, mathtools). The delivered numbering is the unit's own. Assets: no retained span contains a figure environment, a graphics-inclusion command, a PDF-page inclusion, a table or a TikZ picture; no image file is copied into this package, the package declares no asset dependency and no image file is redistributed with it. Licence: version arXiv:2312.02825v3 of this article is distributed under the Creative Commons Attribution 4.0 International licence, as recorded in the arXiv licence metadata for this exact version and shown on the abstract page https://arxiv.org/abs/2312.02825v3. A full rights notice is delivered at licenses/arxiv-2312.02825-CC-BY-4.0.txt. Characters: every retained excerpt is pure ASCII, so the delivered engine, XeLaTeX with its default Latin Modern text font, reports no missing character.
\begin{definition}
    The system's Lagrangian $L \colon TQ \to \mathbb{R}$ is given by
    \begin{align}
        L(\vec{q},\dot{\vec{q}}) = T(\vec{q},\dot{\vec{q}}) - V(\vec{q}) = \frac{1}{2} \dot{\vec{q}} \cdot \vec{M}(\vec{q}) \dot{\vec{q}} -V(\vec{q}), \label{lagrangefunction}
    \end{align}
    where $T:TQ \rightarrow \mathbb{R}$ is the kinetic energy, $\vec{M}(\vec{q}) \in \mathbb{R}^{d \times d}$ denotes the symmetric and positive-definite mass matrix and $V: Q \rightarrow \mathbb{R}$ is a potential function.
\end{definition}

\begin{definition}
    Given a Lagrangian \eqref{lagrangefunction} the quantity
    \begin{align}
        E(\vec{q},\vec{v},\vec{p}) = \vec{p} \cdot \vec{v} - L(\vec{q},\vec{v}) \label{generalized_energy}
    \end{align}
    is called a generalized energy function.
\end{definition}

\begin{subequations}
    \label{HP_EL_discrete}
    \begin{align}
        \vec{q}\npe - \vec{q}\n & = h \, \vec{v}\npeh \label{HP_EL1_discrete} ,                                                           \\
        \vec{p}\npe - \vec{p}\n & = h \DG_1 L(\vec{q},\vec{v})  \label{HP_EL2_discrete} - h\, \sum_{k=1}^{m} \lambda_k \DG g_k(\vec{q}) , \\
        \vec{p}\npeh            & = \DG_2 L(\vec{q},\vec{v}) \label{HP_EL3_discrete} ,                                                    \\
        \vec{g}(\vec{q}\npe)    & = \vec{0} \label{HP_EL4_discrete} ,
    \end{align}
\end{subequations}

\begin{definition}
    A discrete version of the energy function \eqref{generalized_energy} is given by
    \begin{align} \label{discrete_energy_function}
        E\n & = \vec{p}\n \cdot \vec{v}\n - \frac{1}{2}\vec{v}\n \cdot \vec{M}(\vec{q}\n) \vec{v}\n + V(\vec{q}\n) .
    \end{align}
\end{definition}

\begin{proposition}
    The generalized energy \eqref{discrete_energy_function} is conserved by time-stepping method \eqref{HP_EL_discrete} in a discrete sense.
\end{proposition}

\begin{proof}
    Scalar multiplying \eqref{HP_EL2_discrete} with $\Delta \vec{q}$ and adding the scalar product of \eqref{HP_EL3_discrete} with $\Delta \vec{v}$ yields
    \begin{align}
        (\vec{p}\npe -\vec{p}\n) \cdot (\vec{q}\npe - \vec{q}\n) + h \, \vec{p}\npeh \cdot (\vec{v}\npe - \vec{v}\n)                                                           & =                                                                                                                  \\
        h \, \DG_1 L(\vec{q},\vec{v}) \cdot                                                                           (\vec{q}\npe - \vec{q}\n) - h\, \sum_{k=1}^{m} \lambda_k & \DG g_k(\vec{q}) \cdot (\vec{q}\npe - \vec{q}\n) + \DG_2 L(\vec{q},\vec{v}) \cdot (\vec{v}\npe - \vec{v}\n) \notag
    \end{align}
    Inserting \eqref{HP_EL1_discrete} and making use of the directionality property gives
    \begin{align}
        (\vec{p}\npe -\vec{p}\n) \cdot \vec{v}\npeh + \vec{p}\npeh \cdot (\vec{v}\npe - \vec{v}\n) & = L(\vec{q}\npe,\vec{v}\npe) - L(\vec{q}\n,\vec{v}\n)                                          \\
                                                                                                   & \qquad \qquad - \sum_{k=1}^{m} \lambda_k \left(g_k(\vec{q}\npe)-g_k(\vec{q}\n)\right) . \notag
    \end{align}
    After cancelling some terms on the left-hand side of the previous relation and taking into account that the constraints are identically satisfied in each time step (see \eqref{HP_EL4_discrete}), one eventually arrives at
    \begin{align}
        E(\vec{q}\npe,\vec{v}\npe, \vec{p}\npe) = E(\vec{q}\n,\vec{v}\n, \vec{p}\n)
    \end{align}
    which concludes the proof.
\end{proof}

\end{document}
